\documentclass[10pt,journal,letterpaper]{IEEEtran}
\usepackage[margin=1in]{geometry}
\usepackage{cite,amsmath,amssymb,graphicx,booktabs,url}
\usepackage[hidelinks]{hyperref}
\title{Forecasting Interaction Onsets from Aircraft Surveillance with Compact Relational Features}
\author{Junwen Li\thanks{Junwen Li is with Sichuan University, Chengdu, China (e-mail: 2024141460580@stu.scu.edu.cn).}}

% Submission text sizes: retain 10 pt in abstracts, captions, tables and references.
\renewcommand{\small}{\normalsize}
\renewcommand{\footnotesize}{\normalsize}
\raggedbottom
\begin{document}
\maketitle
\begin{abstract}
Forecasting new interactions from aircraft surveillance requires separating future event onsets from proximity already present at prediction time. We formulate an episode-count target that excludes uninterrupted persistence and compare compact aggregate, observed relational, and constant-velocity features under identical data partitions. Low-capacity models predict counts and threshold-exceedance probabilities at horizons of 30, 120, and 300 seconds. Evaluation uses 2,605 scenes from eight held-out source groups at one airport and 1,268 scenes from ten groups at a second airport, without adaptation. Observed pair geometry yields lower mean count error at the development airport, with better-resolved gains at shorter horizons. The transferred system achieves lower mean count and probability errors than training-derived references at 30 and 120 seconds, although relational features do not consistently outperform aggregate features across airports. Constant-velocity projection provides no stable additional validation benefit. These results support a compact surveillance representation for geometric interaction forecasting while delimiting its transferability; the target is not a validated measure of workload or safety risk.
\end{abstract}
\begin{IEEEkeywords}
Aircraft surveillance, event forecasting, relational features, probabilistic prediction, transfer evaluation.
\end{IEEEkeywords}

\section{Introduction}
\IEEEPARstart{M}{ulti-aircraft} surveillance streams contain relational information that traffic counts alone cannot describe. Graph indicators capture evolving air-traffic interactions \cite{isufaj2021spatiotemporal}, while supervised graph models predict controller workload or airspace complexity \cite{pang2023workload,li2024mastgnn}. Public Automatic Dependent Surveillance--Broadcast (ADS-B) archives, however, rarely supply synchronized workload labels. Trajectory benchmarks such as TrajAir and TartanAviation support a complementary task based on observable geometry \cite{patrikar2022trajair,patrikar2025tartanaviation}.

Recent work forecasts relevant aircraft interactions in en-route airspace \cite{henderson2026relevantpairs}, while multimodal trajectory models predict future aircraft motion \cite{prutsch2026ascent}. Here the endpoint is the number of newly arising terminal proximity episodes, with persistence at the cutoff explicitly excluded. If a future-window counter starts with no active pairs, encounters already present at the first future sample are counted as new. Short-horizon performance can then reflect persistence recognition. Separately, a richer representation can appear transferable when it mainly reproduces traffic exposure or when overlapping windows are treated as independent observations.

We formulate surveillance-derived proximity-onset forecasting by initializing the event state at the forecast cutoff, separating future observed onsets from ongoing proximity. Using only reports available at that cutoff, a matched comparison of aggregate state, observed pair geometry, and closest-point-of-approach (CPA) projection assesses the predictive information in each representation. Established low-capacity predictors serve as controlled probes of the onset formulation and representations, keeping estimator capacity comparable across feature sets. Held-out source groups and transfer without airport-specific tuning assess the empirical limits of transportability. Unlike conflict-probability models based on trajectory uncertainty \cite{kuchar2000review}, our probabilities describe exceedance of an observed episode count, not collision risk.

\section{Surveillance Representation and Prediction}
\subsection{Causal Scenes and Interaction Onsets}
We use raw TartanAviation reports for KBTP and KAGC \cite{patrikar2025tartanaviation}. The observation region is a 10-NM-radius cylinder extending from 150~m below to 1,524~m above airport elevation. Each selected source contributes a two-hour window starting at the 10-min bin with the most distinct aircraft. Scenes occur every 10~s and require at least three active aircraft. Invalid identifiers, timestamps, and required states are excluded; duplicate natural keys are resolved deterministically. Each aircraft's latest report is carried forward for at most 30~s.

The input buffer contains twelve causal snapshots from $t-110$~s to cutoff $t$. Features use current states and approximately 30-s changes rather than a learned temporal sequence. Retaining only the final six snapshots reproduces every A/B/C feature on the evaluation scenes. This is input equivalence, not evidence that a shorter history improves prediction.

Let $P(\tau)$ contain observed aircraft pairs whose horizontal and vertical separations are below 5~NM and 1,000~ft. For $\tau_k=t+10k$~s, define
\begin{equation}
Y_h(t)=\sum_{k=1}^{h/10}|P(\tau_k)\setminus P(\tau_{k-1})|,
\quad P(\tau_0)=P(t),
\label{eq:onset}
\end{equation}
where $h\in\{30,120,300\}$~s. A pair active at $t$ contributes nothing while continuously active; clearing and re-entering produces another episode. Thus $Y_h$ counts episodes, not distinct pairs. Future sets use post-cutoff reports with the same staleness bound. Report disappearance and reappearance can also produce observed onsets. Horizontal target distances use the haversine formula with Earth radius 6,371,008.8~m; altitude conversion does not reconcile different physical altitude references.

For example, a pair already close at $t$ and still close at $t+10$~s contributes zero onsets. If it clears at $t+20$~s and returns at $t+30$~s, it contributes one. This definition separates event initiation from occupancy while retaining recurrent episodes; it does not infer unobserved continuous trajectories between surveillance samples.

We predict both $Y_h$ and $\Pr(Y_h\geq q_h\mid x_t)$. The training-only 75th percentile, using the higher-quantile rule, gives $q_h=1,2,5$ for the three horizons. These are statistical event thresholds, not regulatory separation or risk levels.

\subsection{Matched Feature Sets}
Representation A contains eight variables: aircraft count; altitude mean, standard deviation, and range; speed mean and standard deviation; minimum horizontal pair distance; and the fraction below 3,000~m absolute reported altitude. A therefore already includes one pair-distance summary.

Representation B adds five variables: available pair count, active 5-NM/1,000-ft pair count, closing-pair count within 10~NM, mean horizontal pair distance, and minimum vertical pair separation. For each aircraft, the prior observation is the available state 20--45~s earlier closest to a 30-s lag. Closing pairs require both prior states and decreasing horizontal separation. The resulting 13-dimensional vector summarizes both traffic exposure and observed relational structure.

Representation C adds nine CPA variables: projected close-pair count and minimum projected horizontal and vertical separation at each horizon. Feature coordinates are $x=111320(\lambda-\lambda_0)\cos(\pi\phi_0/180)$ and $y=111320(\phi-\phi_0)$ in metres, with longitude and latitude in degrees and the airport as origin. Given relative position $r$ and relative velocity $v$, estimated with each aircraft's actual observation interval,
\begin{equation}
s_h=\operatorname{clip}\left(-\frac{r_{xy}^{\mathsf T}v_{xy}}{\|v_{xy}\|^2},0,h\right).
\label{eq:cpa}
\end{equation}
Both separations are evaluated at $s_h$; it is zero when $\|v_{xy}\|^2\leq10^{-9}$. Pairs missing either prior state are omitted, and empty distance minima use a 100,000-m sentinel. This comparison tests one constant-velocity construction, not all trajectory-based prediction.

\subsection{Estimators and Reference Predictions}
Identical scene IDs, targets, splits, and training-only preprocessing are used across A/B/C. Count candidates are standardized Ridge on $\log(1+Y_h)$ with regularization 1.0, unpenalized Poisson regression, and histogram gradient boosting with Poisson loss. Validation selects Ridge at 30~s and Poisson at 120 and 300~s. Standardized $L_2$-regularized Logistic regression with $C=1$ predicts exceedance probability; seed 1701 is used where applicable. For a single representation shared across horizons, we retain B as the parsimonious specification. B yields lower validation date-equal count error and Brier score than A throughout. C adds nine features and lowers the selected count models' errors at 30/120~s, but has higher count error at 300~s and higher Brier scores at 120/300~s. This development choice favors a compact shared representation; it does not imply horizon-specific optimality. Development scores and model settings appear in the supplementary material.

\begin{table*}[!t]
\caption{Held-out prediction with KBTP-trained models. Values are equal-group means; lower is better. Win counts compare B with the column's reference. The count and probability columns evaluate distinct fitted models.}
\label{tab:results}
\centering
\setlength{\tabcolsep}{4pt}
\begin{tabular}{llrrrrrrr}
\toprule
Airport & $h$ (s) & B MAE & A MAE & Exposure MAE & Wins vs. A/exposure & B Brier & Prevalence & Wins \\
\midrule
KBTP & 30 & 0.580 & 0.592 & 0.630 & 7/8, 8/8 & 0.183 & 0.213 & 7/8 \\
 & 120 & 1.589 & 1.623 & 1.622 & 5/8, 6/8 & 0.200 & 0.239 & 8/8 \\
 & 300 & 3.109 & 3.129 & 3.136 & 6/8, 5/8 & 0.187 & 0.220 & 8/8 \\
\midrule
KAGC & 30 & 0.352 & 0.343 & 0.399 & 4/10, 9/10 & 0.1720 & 0.1728 & 6/10 \\
 & 120 & 1.029 & 1.032 & 1.098 & 6/10, 8/10 & 0.200 & 0.210 & 6/10 \\
 & 300 & 2.041 & 2.024 & 2.116 & 5/10, 7/10 & 0.180 & 0.173 & 7/10 \\
\bottomrule
\end{tabular}
\end{table*}

On the three validation dates, the selected count models give A/B/C mean errors of 0.8029/0.7842/0.7484 at 30~s, 2.1039/2.0205/1.9940 at 120~s, and 4.2415/4.1418/4.1567 at 300~s. Corresponding B/C Brier scores are 0.2086/0.2049, 0.1509/0.1526, and 0.1575/0.1601. These values make the tradeoff explicit: projection helps some validation comparisons but supplies no uniform improvement across the two tasks and three horizons. The shared B specification therefore tests a compact observed representation rather than a horizon-wise combination of feature sets.

The selected Ridge output is $\max\{\exp(\widehat z)-1,0\}$, where $\widehat z$ predicts $\log(1+Y_h)$. Poisson and Logistic fits use at most 1,000 iterations; Logistic uses the lbfgs solver. The boosting candidate uses Poisson loss, 100 iterations, at most 15 leaves, learning rate 0.05, and $L_2$ regularization 1.0. Standardization parameters are estimated only on KBTP training scenes and reused at both evaluation airports.

For scene $i$ with $n_i$ aircraft, the exposure reference predicts $m_i\widehat\lambda_h$, where $m_i=n_i(n_i-1)/2$ and $\widehat\lambda_h=\sum_{j\in\mathrm{train}}Y_{jh}/\sum_{j\in\mathrm{train}}m_j$. The probability reference is the KBTP training exceedance frequency. Both retain their training values at KAGC. Training means, cutoff active-pair counts, and CPA counts provide additional simple count references. These controls separate predictive information from exposure and present-state persistence.



\section{Evaluation and Results}
\subsection{Partitions and Statistical Units}
KBTP development comprises 6,008 scenes on 17 UTC dates: 3,318 training scenes on ten dates, 1,247 validation scenes on three dates, and 1,443 retrospective test scenes on four dates. After completion of model development, eight unused KBTP source groups were selected by a predefined temporal-stratification procedure using authenticated row count only, without inspecting targets or model outputs, yielding 2,605 confirmation scenes. KAGC contributes 1,268 scenes from ten unused groups, one high-volume source per represented month from November 2021 to August 2022. Sources crossing a UTC boundary remain single groups. KAGC supplies no feature selection, normalization, threshold selection, or parameter tuning.

Count mean absolute error (MAE) and probability Brier score are computed within groups and averaged with equal group weight. Overlapping scenes are not treated as independent replicates. Proper probability scores and calibration diagnostics follow \cite{gneiting2007proper,gneiting2007calibration}. The predefined KBTP criterion requires lower MAE than both A and exposure, lower Brier than prevalence, and at least five of eight group wins in each comparison. Transfer requires lower MAE than exposure and lower Brier than prevalence, each with at least six of ten wins. The overall evaluation criterion requires improvement at two or more horizons. These are descriptive criteria, not significance tests.

Paired group differences are bootstrapped 20,000 times (seed 170104) for pointwise percentile 95\% intervals. Two-sided sign tests exclude ties, with Holm correction over 18 airport--horizon--comparison tests. Eight or ten groups limit precision, and calendar separation does not guarantee independent operating conditions.



\subsection{Target Validity and Held-Out Prediction}
Initializing the event counter with $P(t)$ changes the development totals from 18,615, 27,984, and 44,743 episodes to 3,362, 12,731, and 29,490. The largest reduction occurs at 30~s, demonstrating how much a window-initialized target mixes persistence with onset. A separate reconstruction implementation reproduces the onset-target scenes exactly. This verifies target construction rather than forecasting accuracy.

Table~\ref{tab:results} shows that B meets the KBTP criteria at every horizon. Its gains over A are modest. B-minus-A MAE differences are $-0.012$, $-0.034$, and $-0.020$, with respective 95\% intervals $[-0.021,-0.003]$, $[-0.061,-0.007]$, and $[-0.137,0.131]$. The shorter-horizon mean gains are better resolved; the 300-s interval spans zero. No sign test remains significant at 0.05 after Holm correction. The probability comparison shows lower mean Brier scores than unconditional prevalence, not superiority over every alternative representation. Pooled ten-bin expected calibration errors are 0.027, 0.050, and 0.063.

KAGC has lower mean targets than KBTP training: 0.345 versus 0.446, 1.258 versus 1.653, and 2.749 versus 3.739. The 30-s Brier improvement is small (0.1720 versus 0.1728). At 300~s, Brier worsens despite seven group wins, reflecting the magnitudes of the group differences. B outperforms A in mean count error only at 120~s at KAGC. Thus the predefined transfer criteria are met at 30 and 120~s, but a universal advantage of additional relational features is not observed.

\subsection{Sensitivity and Scope}
The supplementary material summarizes target sensitivity, distributional prediction, local-history references, and reproduction settings. Changing the staleness bound to 15 or 60~s changes the 120-s development episode total from 12,731 to 12,601 or 13,471. This does not establish prediction robustness to altered staleness. Separate geometric and exceedance-threshold checks also leave long-horizon probability transfer threshold-sensitive.

An exploratory negative-binomial analysis conducted after the primary evaluation improves count likelihood over matched Poisson models, but its exceedance probabilities do not consistently improve on Logistic regression. Relative to matched Poisson models, its group-equal negative log likelihood is lower by 0.062/0.272/0.891 at KBTP and 0.009/0.131/0.489 at KAGC for 30/120/300~s. However, the joint likelihood-and-probability comparison criterion is met only at KBTP 120~s and KAGC 30/120~s. A fuller count distribution therefore need not improve the specific threshold-exceedance prediction task. The supplementary material gives the comparison rule and probability-score differences.

When past KAGC target labels are available, a chronological local median yields lower 30-s count MAE than B (0.268 versus 0.352) on nine groups after warm-up. This comparison uses additional local label information absent from the primary transfer setting; usefulness depends on the prediction objective and available history. These supplemental results are not additional independent transfer trials.

\subsection{Observation Coverage and Calibration}
The observed-pair and CPA representations depend on the availability of earlier reports. Both aircraft have usable prior states for 18,502 of 20,536 current KBTP pairs (90.1\%) and 3,438 of 5,527 KAGC pairs (62.2\%). Among eligible pairs, relative horizontal velocity is degenerate for 0.07\% and 8.14\%, respectively, setting CPA time to zero. Thus the comparison concerns the CPA features that can be constructed from the available reports; it does not isolate the benefit of ideal motion estimates. These coverage differences do not establish why C lacks consistent gains.

Freshness strata defined by maximum current-report age of at most 10~s versus greater than 10~s contain 1,950/655 KBTP scenes and 329/939 KAGC scenes. KBTP Brier scores are below prevalence in both strata at all horizons. At KAGC, the 30-s Brier scores are 0.149 versus 0.162 for the fresher stratum and 0.178 versus 0.176 for the older stratum. Differences in traffic composition and prevalence prevent a causal interpretation of this association. Reliability summaries use ten equal-width probability bins, and upper KAGC bins can contain only one or two scenes; apparent bin-level deviations are therefore imprecise.

\subsection{Practical Scope}
A local CPU timing sample of 31 scenes with 3--11 aircraft gives warm-call median and 95th-percentile times of 1.68 and 2.43~ms for history-to-feature computation and all three count and three probability outputs. The feature helper also computes unused CPA quantities. These measurements exclude network transfer, disk access, buffer maintenance, and initial history accumulation, so they describe computation at the observed scale rather than end-to-end deployment latency.

The selected high-activity windows do not represent full-day traffic. Missing reports can affect the target, and usable prior-state coverage is lower at KAGC. Low-capacity relational summaries cannot resolve unobserved intent, weather, or controller actions. The target measures observed geometry and requires separate operational validation before it can support workload or safety decisions.

\section{Conclusion}
Cutoff-initialized episode counts distinguish future interaction onsets from uninterrupted proximity. Matched comparisons show that a compact observed-pair representation yields modest reductions in mean prediction error at KBTP, especially at shorter horizons, without stable benefit from the tested CPA extension. The predefined transfer criteria are met at KAGC at 30 and 120~s, but not every representation gain is reproduced. The resulting task and evaluation provide a basis for testing richer surveillance representations without conflating persistence, overlapping samples, or airport-specific adaptation with generalization.

\section*{Data and Code Availability}
The raw surveillance data are available from TartanAviation \cite{patrikar2025tartanaviation} under its applicable terms. Code, derived data, experiment configurations, and result summaries are available in the \href{https://github.com/wenmu663559/adsb-future-onset-interaction-forecasting}{public GitHub repository}. The accompanying supplementary material provides additional analyses, data partitions, and reproduction settings.

\section*{Acknowledgment}
OpenAI Codex assisted with manuscript restructuring and language revision of the existing methods and results.

\clearpage
\bibliographystyle{IEEEtran}
\bibliography{references}
\end{document}
